\section*{Chapter \ref{ch:b2:metamorphic-facies} --- Metamorphic Reactions and Facies}

\begin{solution}{exo:b2:metamorphic-facies:1}
(a) Contact \omterm{def:b2:metamorphic-facies:metamorphism}{metamorphism}, a local thermal \omterm{def:b2:metamorphic-facies:metamorphism}{metamorphism} by the heat of the intrusion.
(b) Dislocation \omterm{def:b2:metamorphic-facies:metamorphism}{metamorphism}, in a fault zone: grain-size reduction, a mylonite.
(c) Ocean-floor \omterm{def:b2:metamorphic-facies:metamorphism}{metamorphism}, by hot seawater near the ridge, with incomplete
recrystallisation. (d) Burial \omterm{def:b2:metamorphic-facies:metamorphism}{metamorphism}, regional, without deformation or
schistosity.
\end{solution}

\begin{solution}{exo:b2:metamorphic-facies:2}
$\Delta_r S = 95.40 - 82.80 = \qty{12.60}{J.K^{-1}}$, $\Delta_r V = 4.986 - 4.415 =
\qty{0.571}{J.bar^{-1}}$; slope $12.60/0.571 = \qty{22.1}{bar/K} = \qty{2.21}{MPa/K}$.
Over \qty{100}{\celsius} the equilibrium pressure rises by \qty{2.21}{kbar}.
\end{solution}

\begin{solution}{exo:b2:metamorphic-facies:3}
(a) Greenschist. (b) Blueschist. (c) Amphibolite: hornblende with plagioclase of more
than \qty{17}{\percent} anorthite. (d) Eclogite.
\end{solution}

\begin{solution}{exo:b2:metamorphic-facies:4}
(a) $\rho = M/V$: kyanite $162.046/44.15 = \qty{3.670}{g.cm^{-3}}$, andalusite
$162.046/51.52 = 3.145$, sillimanite $162.046/49.86 = 3.250$. (b) Kyanite: at higher
pressure the term $V\,(P - 1)$ of the Gibbs energy grows least for the smallest volume.
(c) Natural crystals contain other elements (ferric iron, for instance, replaces some
aluminium in natural andalusite and sillimanite) and the measurements on natural
specimens are less precise than the molar volumes of the tables.
\end{solution}

\begin{solution}{exo:b2:metamorphic-facies:5}
(a) $\Delta_r G = 3900 - 298.15\times8.59 = \qty{+1339}{J} > 0$: kyanite. (b) At
\qty{1}{bar}, $\Delta_r H = T\Delta_r S$: $T = 3900/8.59 = \qty{454}{K} =
\qty{181}{\celsius}$. (c) $P = (8.59\times673.15 - 3900)/0.737 = 1882/0.737 =
\qty{2554}{bar} = \qty{2.55}{kbar}$, against 2.54 measured. (d) The reaction is too slow
at low temperature to proceed: andalusite survives metastably, as most high-temperature
minerals do at the surface.
\end{solution}

\begin{solution}{exo:b2:metamorphic-facies:6}
(a) $\Delta_r H = 7700 + 1000\times(121.39 - 122.29) = \qty{6800}{J}$; $\Delta_r S = 296.44
- 285.02 = \qty{11.42}{J.K^{-1}}$; $\Delta_r G = 6800 - 1000\times11.42 = \qty{-4620}{J}$.
(b) With the values at \qty{298.15}{K}: $7700 - 1000\times12.60 = \qty{-4900}{J}$.
(c) Sillimanite. The simple method errs by \qty{280}{J}, about \qty{6}{\percent}; divided
by $\Delta_r V = 0.571$, that is about \qty{490}{bar} on the line, much less than the
effect of the uncertainty of the enthalpies.
\end{solution}

\begin{solution}{exo:b2:metamorphic-facies:7}
(a) $\Delta_r G = -1127.4 + 1128.5 = \qty{1.1}{kJ}$ and $\Delta_r V = (34.15 -
36.93)/10 = \qty{-0.278}{J.bar^{-1}}$: $P = 1100/0.278 = \qty{3957}{bar} \approx
\qty{4.0}{kbar}$. (b) $\Delta_r S = 88.0 - 91.7 = -3.7$; slope $-3.7/-0.278 =
\qty{13.3}{bar/K}$. (c) With $\Delta_r H = 0$, $P = T\Delta_r S/\Delta_r V = 13.31\,T$:
at \qty{573.15}{K}, \qty{7629}{bar} $= \qty{7.63}{kbar}$. (d) Shell aragonite is
metastable at the surface: organisms make it, and it persists because it turns into
calcite only slowly.
\end{solution}

\begin{solution}{exo:b2:metamorphic-facies:8}
(a) $-32.4/-1.698 = \qty{19.1}{bar/K}$. (b) $P = (T\Delta_r S - \Delta_r H)/\Delta_r V =
(-32.4\times773.15 + 5000)/(-1.698) = (-25050 + 5000)/(-1.698) = \qty{11808}{bar}
\approx \qty{11.8}{kbar}$. (c) $\sqrt{2.6^2 + 3.6^2 + 1.0^2} = \qty{4.55}{kJ}$ on
$\Delta_r H$, so $4550/1.698 = \qty{2.7}{kbar}$ on the pressure. (d) At the pressures of
the eclogite facies the albite component of plagioclase is unstable relative to
jadeite and quartz; the sodium goes into the pyroxene (omphacite), and plagioclase
disappears.
\end{solution}

\begin{solution}{exo:b2:metamorphic-facies:9}
(a) From 1 to \qty{2}{kbar}: $1000/34 = \qty{29.4}{bar/K}$; from 2 to 3: $1000/25 =
\qty{40.0}{bar/K}$. (b) With $\Delta_r S$ fixed, $\Delta_r V$ is inversely proportional to
the slope: it shrank by a factor $40.0/29.4 = 1.36$. (c) The water released: its molar
volume falls as pressure rises, while the volumes of the solids hardly change.
\end{solution}

\begin{solution}{exo:b2:metamorphic-facies:10}
(a) At \qty{873.15}{K}: kyanite $=$ andalusite $(8.59\times873.15 - 3900)/0.737 =
\qty{4885}{bar}$; kyanite $=$ sillimanite $(12.60\times873.15 - 7700)/0.571 =
\qty{5782}{bar}$; andalusite $=$ sillimanite $(4.01\times873.15 - 3800)/(-0.166) =
\qty{1799}{bar}$. (b) At \qty{4885}{bar}, below the kyanite $=$ sillimanite line (5782),
sillimanite is more stable than kyanite, and above the andalusite $=$ sillimanite line
(1799), more stable than andalusite: the point lies in the sillimanite field, and the
kyanite $=$ andalusite equilibrium there is metastable. (c) At \qty{4}{kbar}: between 1799
and \qty{5782}{bar}, sillimanite. At \qty{7}{kbar}: above \qty{5782}{bar} kyanite beats
sillimanite, and above \qty{4885}{bar} it beats andalusite: kyanite.
\end{solution}

\begin{solution}{exo:b2:metamorphic-facies:11}
(a) $z = P/(\rho g) = 3.75\times10^8/(2700\times9.81) = \qty{14.2}{km}$, uncertain by
$0.25\times10^8/26487 = \qty{0.94}{km}$. (b) $(504 - 10)/14.16 =
\qty{35}{\celsius.km^{-1}}$. (c) The three minerals coexist in equilibrium only at the
\omterm{def:b2:metamorphic-facies:invariant}{invariant point}: the path of the rock must pass through it, and the reactions must
stop there. Most rocks pass beside it and keep one or two polymorphs.
\end{solution}

\begin{solution}{exo:b2:metamorphic-facies:12}
(a) X: $T = 10 + 10\times30 = \qty{310}{\celsius}$ (\qty{583.15}{K}); $P = 3000\times
9.81\times30000 = \qty{8.83e8}{Pa} = \qty{8.83}{kbar}$. Y: $T = 10 + 30\times20 =
\qty{610}{\celsius}$ (\qty{883.15}{K}); $P = 2800\times9.81\times20000 = \qty{5.49}{kbar}$.
(b) X: calcite $=$ aragonite at $13.31\times583.15 = \qty{7.76}{kbar} < 8.83$, aragonite;
albite $=$ jadeite + quartz at $(32.4\times583.15 - 5000)/1.698 = \qty{8.18}{kbar} <
8.83$, jadeite + quartz; kyanite $=$ andalusite at $(8.59\times583.15 - 3900)/0.737 =
\qty{1.51}{kbar}$, below the temperature of the triple point: kyanite. Y: calcite $=$
aragonite at \qty{11.76}{kbar} $> 5.49$, calcite; albite $=$ jadeite + quartz at
\qty{13.9}{kbar}, albite; above the triple point, kyanite $=$ sillimanite at
\qty{6.00}{kbar} $> 5.49$ and andalusite $=$ sillimanite at \qty{1.56}{kbar} $< 5.49$:
sillimanite. (c) X, at \qty{10}{\celsius\per\km}, is the colder path: jadeite and
aragonite at only \qty{310}{\celsius} place it in the blueschist facies. Y, three times
warmer, ends with sillimanite at \qty{610}{\celsius} and moderate pressure, in the
amphibolite facies.
\end{solution}

\begin{solution}{pb:b2:metamorphic-facies:1}
\noindent\textbf{Part I.}

\textbf{1.} $\Delta_r H = -2589.9 + 2593.8 = \qty{3.9}{kJ}$; $\Delta_r S = 91.39 - 82.80 =
\qty{8.59}{J.K^{-1}}$; $\Delta_r V = 5.152 - 4.415 = \qty{0.737}{J.bar^{-1}}$.

\textbf{2.} $\qty{7.7}{kJ}$; $\qty{12.60}{J.K^{-1}}$; $\qty{0.571}{J.bar^{-1}}$.

\textbf{3.} $\qty{3.8}{kJ}$; $\qty{4.01}{J.K^{-1}}$; $\qty{-0.166}{J.bar^{-1}}$; each is
the second minus the first, since andalusite cancels.

\textbf{4.} Sillimanite has the largest entropy (95.40): it lies on the high-temperature
side of both its lines. Kyanite has the smallest volume (4.415): it lies on the
high-pressure side of both its lines. Andalusite, the largest volume, holds the
low-pressure side.

\textbf{5.} $8.59/0.737 = \qty{11.66}{bar/K}$; $12.60/0.571 = \qty{22.07}{bar/K}$;
$4.01/(-0.166) = \qty{-24.16}{bar/K}$.

\textbf{6.} Sillimanite has more entropy but less volume than andalusite: $\Delta_r S > 0$
and $\Delta_r V < 0$. Raising the temperature favours sillimanite, raising the pressure
too, so the line runs to lower pressure at higher temperature.

\noindent\textbf{Part II.}

\textbf{7.} $3900 - 298.15\times8.59 = \qty{+1339}{J}$; $7700 - 298.15\times12.60 =
\qty{+3943}{J}$. Both positive: kyanite is stable at the surface.

\textbf{8.} $T = \Delta_r H/\Delta_r S$: $3900/8.59 = \qty{454.0}{K} = \qty{181}{\celsius}$;
$7700/12.60 = \qty{611.1}{K} = \qty{338}{\celsius}$; $3800/4.01 = \qty{947.6}{K} =
\qty{674}{\celsius}$.

\textbf{9.} Heating at \qty{1}{bar}, kyanite gives way to andalusite at
\qty{181}{\celsius}; at \qty{338}{\celsius} kyanite is no longer stable, so kyanite $=$
sillimanite is metastable there; andalusite gives way to sillimanite at
\qty{674}{\celsius}. The first and third are stable equilibria.

\textbf{10.} $(8.59\times673.15 - 3900)/0.737 = \qty{2554}{bar} = \qty{2.55}{kbar}$,
\qty{0.01}{kbar} from experiment.

\textbf{11.} $(12.60\times973.15 - 7700)/0.571 = 4562/0.571 = \qty{7989}{bar} =
\qty{7.99}{kbar}$, \qty{0.13}{kbar} above experiment.

\textbf{12.} Kyanite $=$ andalusite rises from \qty{181}{\celsius} at \qty{1}{bar} with
slope \qty{11.7}{bar/K}; andalusite $=$ sillimanite falls from about \qty{11.5}{kbar} at
\qty{200}{\celsius} to \qty{1}{bar} at \qty{674}{\celsius}; kyanite $=$ sillimanite, the
steepest, crosses both near \qty{514}{\celsius} and \qty{3.9}{kbar}, each line stable on
one side of the crossing only (the diagram of the chapter).

\noindent\textbf{Part III.}

\textbf{13.} $a_1 = 8.59/0.737 = \qty{11.655}{bar/K}$, $b_1 = 3900/0.737 =
\qty{5291.7}{bar}$; $a_2 = 12.60/0.571 = \qty{22.067}{bar/K}$, $b_2 = 7700/0.571 =
\qty{13485.1}{bar}$.

\textbf{14.} $T = (b_2 - b_1)/(a_2 - a_1) = 8193.4/10.412 = \qty{786.92}{K} =
\qty{513.8}{\celsius}$.

\textbf{15.} $P = 11.655\times786.92 - 5291.7 = \qty{3880}{bar} = \qty{3.88}{kbar}$.

\textbf{16.} $(4.01\times786.92 - 3800)/(-0.166) = (3155.5 - 3800)/(-0.166) =
\qty{3882}{bar}$: the same pressure to within the rounding of the data, as the proposition of the chapter requires.

\textbf{17.} \qty{514}{\celsius} lies \qty{10}{\celsius} from \qty{504}{\celsius}, inside
the \qty{20}{\celsius} uncertainty; \qty{3.88}{kbar} lies \qty{0.13}{kbar} from 3.75,
inside 0.25. The agreement is good.

\textbf{18.} No: \qty{536}{\celsius} lies \qty{32}{\celsius} from experiment, outside its
uncertainty. Neglecting the heat capacities was not the main error: the enthalpies are.

\noindent\textbf{Part IV.}

\textbf{19.} $\Delta_r H$ of kyanite $\to$ andalusite rises by \qty{1000}{J}; at fixed
$T$, $P = (T\Delta_r S - \Delta_r H)/\Delta_r V$ falls by $1000/0.737 = \qty{1357}{bar}$.

\textbf{20.} $\Delta T = -1357/(a_2 - a_1) = -1357/10.41 = \qty{-130}{K}$: the triple point
moves to about \qty{384}{\celsius}.

\textbf{21.} The kyanite $=$ sillimanite line falls by $1000/0.571 = \qty{1751}{bar}$, and
the triple point moves by $+1751/10.41 = \qty{+168}{K}$.

\textbf{22.} Each enthalpy alone, varied by $\pm\qty{2.0}{kJ}$, moves the triple point by
$\pm\qty{260}{K}$ (andalusite) or $\pm\qty{336}{K}$ (sillimanite): calorimetry alone allows
triple points over several hundred degrees.

\textbf{23.} The tabulated enthalpies are not purely calorimetric: within their
uncertainties they are adjusted so that the computed lines fit the reversed experiments.
The experiments fix the diagram; the tables carry it, consistently, to other reactions.

\textbf{24.} The triple point computed from the tables with constant properties lies at
\textbf{\qty[text-series-to-math=true]{514}{\celsius}}.
\end{solution}
